% ============================================================
%  11_visualisation.tex — Section 11 : Visualisation et reporting
% ============================================================
\section{Visualisation et reporting interactif}

L'analyse manuelle des \texttt{data.json} étant fastidieuse, trois scripts produisent des
pages HTML statiques (autonomes, ouvrables en \texttt{file://} dans Chrome) qui permettent
d'explorer les résultats interactivement sans serveur web ni dépendance externe.

\subsection{\texttt{view.py} — viewer par taille}

Deux modes :
\begin{itemize}
  \item \textbf{Mode config} : \texttt{python view.py output/hX/config\_name} scanne le
    dossier et produit le \texttt{data.json} de cette config.
  \item \textbf{Mode agrégé} : \texttt{python view.py output/hX -{}-aggregate} consomme tous
    les \texttt{data.json} des sous-dossiers et produit \texttt{output/hX/view.html}
    embarquant les 8 configs (aucun \texttt{fetch}, aucune dépendance CDN).
\end{itemize}

Le viewer agrégé propose :
\begin{itemize}
  \item \textbf{Sélection de config} en haut (8 boutons), avec mode \textbf{comparaison}
    multi-sélection (panneaux côte-à-côte avec diff-highlight des molécules dont la
    classification diffère entre configs).
  \item \textbf{Cards de stats} : Molécules / Originaux plans / Solutions CSP / Solutions
    stables planes / Instables / Stables non planes (basculées selon présence du bloc
    \texttt{runs}).
  \item \textbf{Toolbar} : recherche par nom, filtres (Originaux, Solutions, Stabilité),
    expand all / collapse all.
  \item \textbf{Légende stabilité} repliable, expliquant les 6 classes et leurs critères.
  \item \textbf{Tableau des molécules} cliquables-dépliables. Pour chaque solution
    multi-runs : badge classification (couleur), \texttt{planar\_count/n}, $\mu \pm \sigma$,
    \textbf{barre de dispersion SVG} (échelle fixe 0--30°, ligne pointillée au seuil 10°,
    ticks pour chaque run individuel, marqueur $\mu$), valeurs \texttt{min--max}.
  \item \textbf{Section ``Couverture CSP''} repliable : analyse de la partition des solutions
    planaires entre les 8 configs (cf. ci-dessous).
  \item \textbf{Bandeau \texttt{batch\_meta}} : si \texttt{output/hX/batch\_meta.json}
    existe, affiche en haut les stats du dernier run \texttt{batch\_all} (durée, nombre
    d'instances, options).
\end{itemize}

\subsection{Section ``Couverture CSP''}

Permet de répondre à la question : \emph{quelles solutions chaque config trouve, lesquelles
sont communes, lesquelles sont uniques à un sous-ensemble} ? Deux visualisations côte-à-côte :

\begin{itemize}
  \item \textbf{Répartition par config (toutes solutions)} : un bar chart horizontal, une
    ligne par config, segmenté par classe de stabilité (mêmes couleurs que les badges).
    Permet de voir d'un coup d'\oe{}il quelle config trouve combien de solutions et leur
    profil qualité.
  \item \textbf{Liste d'intersections} (filtrée aux solutions \texttt{always\_planar} +
    \texttt{mostly\_planar}) : 1 section dépliable par sous-ensemble de configs, triée du
    plus universel (toutes les 8 configs) au plus unique (1 seule config). Cliquer dévoile
    la liste des couples \texttt{(molécule, sizes)} de cette intersection.
\end{itemize}

Identité d'une solution : couple \texttt{(mol\_name, sizes)}, indépendant de la config qui
l'a trouvée. Une solution est dite \emph{universelle} si elle apparaît dans les 8 configs,
\emph{unique} si elle n'apparaît que dans une seule.

\subsection{\texttt{update\_report.py} — rapport global}

Génère \texttt{csp\_solver/experiments/report/index.html}, un rapport unique couvrant toutes
les tailles \texttt{hX} testées. Contenu :

\begin{itemize}
  \item Navigation sticky avec \texttt{IntersectionObserver} pour le suivi de section active.
  \item Sections numérotées : Introduction, Méthodologie (pipeline + xTB + ACP), Résultats,
    Stabilité inter-runs, Observations, Détails, Commandes.
  \item \textbf{Charts SVG} générés en Python pur (pas de dépendance D3/Plotly) :
    \begin{itemize}
      \item Bar chart empilé : solutions planes vs non-planes par taille
      \item Donut chart : pourcentage global de solutions planes
      \item Pie chart : répartition des classes de stabilité
      \item Scatter plot $\mu \times \sigma$ : 1 point par solution, couleur par classe,
        tooltips avec $h$/config/molécule, ligne verticale au seuil 10°
    \end{itemize}
  \item \textbf{Tableau résultats par h} : nombre de molécules, solutions, planes, non
    planes, \% plan (avec barre de progression colorée), originaux plans.
  \item \textbf{Tableau stabilité par h} : 6 colonnes (1 par classe), affiché seulement si
    au moins une donnée multi-runs est disponible.
  \item Lignes dépliables : cliquer une ligne du tableau résultats déplie une mini-table par
    molécule.
\end{itemize}

\subsection{\texttt{bundle.py} — archive portable}

Crée \texttt{experiments\_bundle.zip}, archive autonome contenant :
\begin{itemize}
  \item \texttt{report/index.html} (rapport global)
  \item \texttt{output/hX/view.html} (viewer par taille)
  \item Les fichiers \texttt{*\_opt.xyz} référencés par les liens HTML (anciens single-run et
    nouveaux multi-runs, captés par 3 patrons de glob distincts)
  \item \texttt{README.txt} avec instructions d'ouverture
\end{itemize}

L'archive est destinée à être envoyée à un collaborateur. Le destinataire l'extrait et ouvre
\texttt{report/index.html} dans Chrome — toutes les visualisations fonctionnent en
\texttt{file://}, sans serveur ni installation.

\subsection{Refonte du dashboard : passage en MD-only}

\dateblock{1\textsuperscript{er} mai 2026}

Le dashboard \texttt{report/index.html} a été \textbf{repensé en true dashboard} :
information cruciale visible immédiatement, sections secondaires repliées par défaut.
Suite à la convergence MR/MD à 98\%, le choix a été fait de \textbf{ne plus mentionner
les multi-runs dans le résumé} et de présenter uniquement le verdict MD.

\textbf{Hero KPIs (4 tuiles uniformes en haut de page).} Format régulier :
chiffre absolu en gros (BIG), pourcentage ou contexte en petit (SUB). Plus de
mélange chiffres absolus / pourcentages dans la même rangée.

\begin{center}
\begin{tabular}{lll}
  \toprule
  Tuile & BIG (compte) & SUB (\%/contexte) \\
  \midrule
  Molécules     & 25  & h3 à h5 \\
  Solutions CSP & 53  & toutes configurations \\
  \bucketplane{} Plans (MD) & 39  & 74\% des solutions \\
  \bucketnonplane{} Non plans (MD) & 14  & 26\% des solutions \\
  \bottomrule
\end{tabular}
\end{center}

\textbf{Sections masquées du dashboard (mais conservées dans le code).}
Les helpers Python (\texttt{\_compute\_stability}, \texttt{\_render\_alerts},
\texttt{\_render\_divergences\_block}, \texttt{\_render\_md\_section},
\texttt{\_compute\_dashboard\_kpis}) restent appelables, mais
\texttt{generate\_html} ne les rend plus pour le moment :

\begin{itemize}
  \item Section ``Stabilité multi-runs'' (pie + scatter $\mu \times \sigma$) : retirée.
  \item Strip d'alertes (divergences MR/MD) : retiré.
  \item Bloc tableau divergences franches : retiré.
  \item Section ``Validation MD'' redondante avec le hero : retirée.
\end{itemize}

\textbf{Bar chart et tableau ``Par taille''.} Refondus pour afficher 2 segments
(plan/non plan MD) au lieu de 6 classes de stabilité multi-runs. La fonction interne
\texttt{generate\_bar\_chart\_svg} accepte toujours un dict \texttt{\{PLANE, AUTRES, NON\_PLANE\}},
mais \texttt{AUTRES} est forcé à 0 par \texttt{\_solution\_bucket} qui retourne
uniquement \texttt{PLANE} ou \texttt{NON\_PLANE} basé sur \texttt{md\_validation.planar}.
La logique 3-bucket préexistante reste accessible via \texttt{\_solution\_bucket\_3bucket}
si on veut la réactiver.

\textbf{Méthodologie et Commandes.} Repliées par défaut (\texttt{<details>}). La section
\textit{Méthodologie} a été enrichie d'un détail étape par étape (Étapes 1 à 5, du
\texttt{.graph} à l'ACP) avec les paramètres MD complets. La section \textit{Commandes}
classe les commandes en 5 sous-sections, MD en premier (recommandé), puis multi-runs,
puis 1-run, puis régénération des rapports, puis explication des combinaisons de flags CSP.

\subsection{Mini-tableau MD par config dans \texttt{view.html}}

\dateblock{1\textsuperscript{er} mai 2026}

En haut de chaque page \texttt{output/hX/view.html}, juste avant la barre de sélection de
config, un \textbf{petit tableau récap MD} donne d'un coup d'\oe{}il les pourcentages plans
et non-plans par configuration :

\begin{center}
\begin{tabular}{lcccc}
  \toprule
                     & default       & no-freeze     & adj-57       & \ldots \\
  \midrule
  \bucketplane{} Plans (MD)    & 100\% (5/5) & 100\% (5/5) & 100\% (5/5) & \ldots \\
  \bucketnonplane{} Non plans (MD) & 0\% (0/5)   & 0\% (0/5)   & 0\% (0/5)   & \ldots \\
  \bottomrule
\end{tabular}
\end{center}

Lignes colorées (vert pour plan, rouge pour non plan) avec valeur principale en pourcentage
et compte exact \texttt{n/total} en petit dessous. Calculé en JS dans \texttt{view.js}
(fonction \texttt{renderMDSummaryTable}), affiché uniquement si au moins une solution a un
bloc \texttt{md\_validation}.

\subsection{Affichage ``Aucune solution CSP''}

\dateblock{mai 2026 (mi-mois)}

Avec l'exclusion par défaut de la solution tout-hexagones et la régénération de la
table de voisinage par MD, certaines molécules (en particulier les benzénoïdes
entièrement gelés où chaque hexagone a $\deg=6$) se retrouvent avec \textbf{zéro
solution CSP}. Le viewer doit donc distinguer visuellement ces cas.

Trois éléments d'UI ajoutés dans \texttt{view.html} :

\begin{itemize}
  \item Sur la \textbf{ligne molécule} : la cellule \emph{Verdict} affiche une pill
    grise pointillée \texttt{$\oslash$ Aucune solution CSP} (au lieu d'un tiret).
    La cellule \emph{Solutions} montre \texttt{0} dans une étiquette grise. La
    ligne entière reçoit la classe CSS \texttt{.no-sol-row} qui grise le fond et
    italique-grise le nom — la molécule reste lisible mais clairement secondaire.
  \item Sur les \textbf{cards en haut} : si au moins une molécule est dans ce cas,
    une nouvelle card grise \texttt{$\oslash$ Sans solution CSP : N} apparaît
    automatiquement entre \emph{Solutions CSP} et \emph{Plans}. Cachée sinon
    (zéro pollution visuelle quand toutes les molécules ont des solutions).
  \item En \textbf{mode comparaison} (panneaux côte-à-côte) : même traitement,
    version compacte \texttt{$\oslash$ Aucune sol.} pour tenir dans des panneaux
    étroits.
\end{itemize}

Palette : \texttt{\#6a737d} / \texttt{\#eaecef} (gris neutre, distinct des verts /
oranges / rouges des verdicts pour ne pas créer d'ambiguïté).

\subsection{Section ``Couverture CSP'' alignée sur MD}

\dateblock{1\textsuperscript{er} mai 2026}

Le critère d'inclusion d'une solution dans la partition CSP était auparavant :
\textit{always\_planar + mostly\_planar} (multi-runs) avec fallback \texttt{sol.planar}.
Il devient simplement : \texttt{md\_validation.planar = true} (fallback \texttt{sol.planar}).

\textbf{Conséquence sur le bar chart ``Répartition par config''} : il n'affiche plus 6 classes
de stabilité mais 2 segments par ligne (vert plan / rouge non-plan), cohérent avec le
verdict MD. Les fonctions \texttt{computePerConfigBreakdown} et \texttt{renderPerConfigBarSVG}
ont été simplifiées en conséquence.

\subsection{Architecture des templates}

Les templates HTML/CSS/JS sont externalisés dans \texttt{templates/} pour bénéficier de la
coloration syntaxique de l'IDE et permettre l'édition propre. La substitution est faite côté
Python avec \texttt{string.Template} (syntaxe \texttt{\$variable}).

\begin{lstlisting}[style=benzai]
csp_solver/experiments/templates/
  common.css        # Palette + classes partagees (cards, planar/non-planar)
  view.html         # Template du viewer (substitution $h_name, $configs_json, ...)
  view.css          # Styles specifiques au viewer
  view.js           # Logique JS du viewer (IIFE avec etat prive,
                    # fonctions onclick exposees sur window pour le HTML inline)
  report.html       # Template du rapport global
  report.css        # Styles specifiques au rapport
  report.js         # Smooth scroll + IntersectionObserver
\end{lstlisting}

\textbf{Décisions architecturales} :
\begin{itemize}
  \item Le HTML est statique : toutes les données sont injectées au moment de la
    génération via \texttt{<script>window.\_\_DATA\_\_ = \{\ldots\};</script>}. Aucun fetch,
    aucune dépendance CDN, aucun build step.
  \item Le calcul interactif (filtres, tri, partition CSP, classification) est intégralement
    fait \textbf{côté JS} (à l'exception des charts du rapport global, générés en Python).
    Cela évite d'élargir le payload \texttt{data.json}.
  \item Les helpers SVG côté Python (\texttt{generate\_bar\_chart\_svg},
    \texttt{generate\_donut\_svg}, \texttt{generate\_pie\_svg},
    \texttt{generate\_scatter\_svg}) sont des fonctions pures retournant une chaîne
    \texttt{<svg>...</svg>}, faciles à tester et à réutiliser.
\end{itemize}
